% Vorgaben Assignment aus Studienheft SQL03
% Formatvorgaben fuer den Text
% Umfang: 8 - 10 Seiten (inkl. Abbildungen und Tabellen, aber ohne Deckblatt, % Gliederung und Literaturverzeichnis, Eidesstattliche Erklaerung)
% Zeilenabstand: 1,5
% Schriftart: frei
% Schriftgrad: 12 pt
% Variablen, physikalische Groessen und Funktionszeichen werden kursiv gedruckt.
% Korrekturrand: links: 4,5 cm, rechts 2,0 cm, oben und unten jeweils 3,0 cm
% Deckblatt: (Adresse, AKAD-E-Mail-Adresse, Immatrikulationsnummer, Modul-
% bezeichnung, Thema, Datum, Felder für Korrektor)
% Gliederung (1 Seite)
% Literaturverzeichnis (3 - 5 Literaturquellen  z. B. Lehrbuecher, aktuelle Fachartikel recherchieren)
% Eidesstattliche Erklaerung (unterschrieben und fest eingebunden)
% Bearbeitungsdauer: 2 Monate


\documentclass[a4paper,12pt]{article}
\usepackage[ngerman]{babel}
\usepackage[nottoc]{tocbibind} % Anzeigen des Literaturverzeichnisses im TOC
\usepackage{epsfig}
\usepackage{times}
\usepackage{supertabular}
\usepackage{wrapfig}
\usepackage{multirow}
\usepackage[onehalfspacing]{setspace}
\usepackage{listings}
\usepackage{mathptmx}
\usepackage{geometry}
\usepackage{helvet}
\usepackage{courier}
\usepackage{setspace}
\usepackage{textcomp}
\usepackage[T1]{fontenc}
\usepackage[utf8]{inputenc}
\usepackage{fancyhdr}
\usepackage{float} % Notwendig fuer figure[h]
\usepackage[printonlyused]{acronym}

\newif\iflistoffigures
\newif\iflistoftables
\newif\ifacronym

%Titel
\newcommand*{\Titel}{SWE02 -- Sicherheitsaspekte eines Webshops} 

%Betreff
\newcommand*{\Betreff}{Assignment im Modul SWE02 -- \\ Softwareentwicklung \LaTeX{}} 

%Betreuer
\newcommand*{\Betreuer}{Betreuer: Prof. Kirchberg} 

%Vor- und Nachname
\newcommand*{\Name}{Stefan Waidele}

%Straße und Hausnummer
\newcommand*{\Strasse}{Ensisheimer Straße 2} 

%Plz und Ort
\newcommand*{\PlzOrt}{79395 Neuenburg am Rhein} 

%Immatrikulationsnummer
\newcommand*{\Immatrikulationsnummer}{102 81 71}

%Email 
\newcommand*{\Email}{Stefan.Waidele@AKAD.de} 


% Verzeichnisse (Wenn nicht benötigt, Zeile mit % auskommentieren oder löschen

%% Abbildungsverzeichnis 
%\listoffigurestrue
%% Tabellenverzeichnis
%\listoftablestrue
%% Abkürzungsverzeichnis
\acronymtrue


% Wittwen und Waisen verhindern
\clubpenalty10000
\widowpenalty10000
\displaywidowpenalty=10000

\usepackage[flushmargin,hang,ragged]{footmisc}
\usepackage{lmodern} %Type1-Schriftart für nicht-englische Texte
\usepackage{fancyhdr}

\usepackage[
	pdftitle={\Titel},
	pdfsubject={SWE02 -- Sicherheitsaspekte einer Web--Anwendung},
	pdfauthor={Stefan Waidele},
	pdfkeywords={akad, swe02, assignment, wirtschaftsinformatik}
	hyperfootnotes=false,
	colorlinks=true,
	linkcolor=black,
	urlcolor=black,
	citecolor=black
]{hyperref}

%\renewcommand{\familydefault}{\rmdefault}
\renewcommand{\bflabel}[1]{\normalfont{\normalsize{#1}}\hfill}

%% Definition for Codeschnipsel im Fließtext
\newcommand{\code}{\texttt}

%% Für Codeblöcke mit Syntax-Highlighting
%% http://www.ctan.org/tex-archive/macros/latex/contrib/minted/
\usepackage{minted}
\definecolor{bg}{rgb}{0.95,0.95,0.95}


\makeatother

\geometry{a4paper, left=45mm, right=20mm, top=30mm, bottom=30mm}
\pagenumbering{roman}

\begin{document}
	
\parskip=1em
\parindent=0cm

%description: Deckblatt in Deutsch
%% Basierend auf einer TeXnicCenter-Vorlage von Tino Weinkauf.
%% sowie der akad-vorlage von Daniel Falkner
%%%%%%%%%%%%%%%%%%%%%%%%%%%%%%%%%%%%%%%%%%%%%%%%%%%%%%%%%%%%%%

%%%%%%%%%%%%%%%%%%%%%%%%%%%%%%%%%%%%%%%%%%%%%%%%%%%%%%%%%%%%%
%% Deckblatt
%%%%%%%%%%%%%%%%%%%%%%%%%%%%%%%%%%%%%%%%%%%%%%%%%%%%%%%%%%%%%
%%
%% ACHTUNG: Sie benötigen ein Hauptdokument, um diese Datei
%%          benutzen zu können. Verwenden Sie im Hauptdokument
%%          den Befehl "\input{dateiname}", um diese
%%          Datei einzubinden.
%%

\begin{titlepage}

%\thispagestyle{empty}

\vspace{5cm}

\Name \\ 
\Strasse \\ 
\PlzOrt\\ 
\href{mailto:\Email}{\Email}

AKAD University\\
Immatrikulationsnummer: \Immatrikulationsnummer

\vfill

\begin{tabbing}
Modul SWE02 --- \=Softentwicklung\\ 
                \>Assignment  
\end{tabbing}
\LARGE
\textsc{Sicherheit bei Web--Anwendungen}\\

\vfill

\normalsize

Betreuer: Prof. Paul Kirchberg

\today %%Datum der Abgabe - am besten selbst reinschreiben.

\vfill

\includegraphics[width=3cm]{akad_logo.png}

AKAD University

\end{titlepage}


%\includegraphics[scale=0.35]{akad_logo.png}

%\clearpage
\normalsize


\normalsize

\begin{spacing}{1.0} % Verzeichnisse werden mit einzeiligem Abstand gesetzt
\parskip=0em
\newpage

% Inhaltsverzeichnis
\setcounter{tocdepth}{2}
\tableofcontents 
\newpage

% Abbildungsverzeichnis
\iflistoffigures
\listoffigures 
\newpage
\fi

% Tabellenverzeichnis
\iflistoftables
\listoftables 
\newpage
\fi

% Abkürzungsverzeichnis
\ifacronym
\section*{Abkürzungsverzeichnis}
\addcontentsline{toc}{section}{Abkürzungsverzeichnis} 
\begin{acronym}[\hspace{3cm}]
	\acro{CAPTCHA}{Completely Automated Public Tourint Test to Tell Computers and Humans Apart}
	\acro{CSS}{Cascading Style Sheets}
	\acro{DoS}{Denial of Service}
	\acro{DDoS}{Distributed Denial of Service}
	\acro{ERP}{Enterprise Resource Planning}
	\acro{HTTP}{Hypertext Tranfsfer Protocol}
	\acro{HTTPS}{Hypertext Tranfsfer Protocol Secure}
	\acro{OWASP}{Open Web Application Security Project}
	\acro{SCM}{Supply Chain Management}
	\acro{SQL}{Structured Query Language}
	\acro{WWW}{World Wide Web}
	\acro{XSS}{Cross Site Scripting}
%%
%	\acro{}{}
%	\acro{}{}
%	\acro{}{}
%	\acro{}{}
%	\acro{}{}
%	\acro{}{}
%	\acro{}{}
%	\acro{}{}
%	\acro{}{}
	\acro{Test}{Wird nicht im Text verwendet und taucht auch nicht im Verzeichnis auf}
\end{acronym}

\fi

\parskip=1em
\end{spacing} 

\clearpage

\newcounter{romanPagenumber} 
\setcounter{romanPagenumber}{\value{page}} % Roemische Seitenanzahl speichern.


\pagestyle{fancy}
\fancyhead{}
\fancyhead[LO,RE]{\textsc{\Titel}}
\fancyhead[RO,LE]{\thepage}
\fancyfoot[CO,CE]{}

\nocite{*} 

\pagenumbering{arabic}

\begin{spacing}{1.5} % Zeilenabstand: 1,5 fuer den Textteil

\section{Einleitung}
\label{sec:einleitung}


\subsection{Ziele dieser Arbeit}

In dierser Artbeit sollten mögliche Angriffstechniken auf Webanwendungen 
sowie geeignete Gegenmaßnahmen dargestellt werden.

Auch wenn eine Internetanwendung wie etwa ein Web--Shop nur einem beschränkten Personenkreis zur Verfügung stehen soll,
so ist zumindest das Anmeldeformular\footnote{Zumeist gilt dies für noch mehr Seiten mit Benutzerinteraktion, wie z.B. die Grundfunktionalität des Warenkorbs, Suchfunktione, etc.} aufgrund der Struktur des \ac{WWW} doch für Angriffe beliebiger Personen erreichbar.
Daher sind schon bei der Implementierung eine Reihe von Sicherheitsaspekten zu beachten, um dieser umfangreichen Bedrohung entgegen zu wirken.

\subsection{Aufbau der Arbeit}

Zunächst werden im Kapitel~\ref{sec:grundlagen} erörtert, warum der Sicherheitsaspekt bei der Erstellung eines Web–Shops ein wichtig ist. Es werden Angriffsziele, Szenarien genannt, sowie die möglichen wirtschaftlichen Konsequenzen eines erfolgreichen Angriffs aufgezeigt. 

Im Kapitel~\ref{sec:hauptteil} wird gezeigt, wie Angriffe mit dem Methoden Brute Force, Man in the Middle, SQL--Injection sowie \ac{XSS} durchgeführt und abgewehrt werden können.

\subsection{Abgrenzung}

Angriffe im Internet können an vielen Punkten der Datenübertragung statt finden. Diese Arbeit beschäftigt sich hauptsächlich mit den Methoden, auf die schon bei der Softwareentwicklung geachtet werden sollte. Risiken, die in den Verantwortungsbereich des Serverbetreibers oder WLAN--Anbieters, oder gar des Kunden (Datendiebstahl bzw. Keylogger auf dem Kunden--PC) liegen, werden hier nicht weiter behandelt, auch wenn sie im Grundlagenteil der Vollständigkeit halber erwähnt werden.

Außerdem wird davon ausgegangen, dass der Webshop den Kaufprozess lediglich bis zur Bestellung abbildet. Weitere Schritte, speziell die Bezahlung, werden auf herkömmlichem Weg abgewickelt.
\section{Grundlagen}
\label{sec:grundlagen}

\subsection{Schützenswerte Güter}

Vertraulichkeit, Integrität, Verbindlichkeit, Verfügbarkeit

\subsection{Angriffsziele}

Webshop, unsere Daten, die Daten unserer Kunden

\subsection{Angriffsarten}

Denial of Service \ac{DOS} bzw. \ac{DDOS}, Man in the Middle, Cross Site Scripting \ac{XSS}, SQL--Injection \ac{SQL}, Brute Force
\section{Angriffsarten und Gegenmaßnahmen}
\label{sec:hauptteil}

Je nach Angriffsziel und Absicht des Angreifers kommen unterschiedlichen Angriffsmethoden zum Einsatz. Das \ac{OWASP} erstellt regelmäßig eine Liste der der wichtigsten Angriffsarten auf Internetapplikationen. Die in dieser Arbeit untersuchten Risiken Injection (Am Beispiel \ac{SQL}) und \acs{XSS}\footnote{X für „Cross“, um Verwechslungen mit \ac{CSS} zu vermeiden.} belegen auf dieser Liste für 2013 die Plätze eins und drei\footnote{vgl. \cite{owasp}}. Die beiden weiteren hier besprochenen Angriffsmethoden \ac{DoS} bzw. \ac{DDoS}, Man in the Middle und Brute Force können als zumindest als Teilsspekte der ebenfalls in der Liste vorkommenden Punkte „Broken Authentication and Session Management“ (Platz 2) und „Security Misconfiguration“ (Platz 5) betrachtet werden.

\subsection{Denial of Service}
\label{sec:dos}

Unter einem \ac{DoS} versteht man einen Angriff auf die Verfügbarkeit eines Angebots. Im Kontext dieser Arbeit verfolgt der Angreifer somit das Ziel, den Webshop für die Kunden unbenutzbar zu machen. Dies kann aufgrund der öffentlichen Erreichbarkeit mit legitim erscheinenden Anfragen an den Webserver geschehen. Aufgrund der Leistungsfähigkeit moderner Webserver ist dies aber i.d.R. nicht mit einzelnen Rechnern zu leisten. Daher werden meist große Netzwerke von Rechnern\footnote{Die Besitzer dieser Rechner wissen i.d.R. nichts von diesen Angriffen, da die Schadsoftware unbemerkt über Computerviren eingeschleust wurde.} für einen gemeinsamen, dann \ac{DDoS} gennanten Angriff verwendet.\footnote{vgl. \cite{carr}, Kapitel 5}

Oft sind \ac{DoS}--Angriffe nicht zweifelsfrei nachweisen, da Server auch durch eine entsprechende Menge von legitimen Zugriffen\footnote{z.B. durch Berichterstattung in populären Medien bzw. auf reichweitenstarken Webseiten („Slashdot--Effekt“)} überlastet werden können. Mögliche Gegenmaßnahmen sind die Nutzung einer skalierbaren Infrastruktur, die allgeimein effektive Implementierung der Shopsoftware sowie eine entsprechende Konfiguration der Serversoftware (z.B. Caching). 

Da \ac{DoS}--Angriffe nicht nur per \ac{HTTP} auf die Shopsoftware selbst sondern auch auf Firewalls oder Router zielen können\footnote{vgl. \cite{amoroso}, S. 60ff} sind die Abwehrmöglichkeiten für den Entwickler der Shopsoftware sehr beschränkt.

Außerdem besteht ein Zusammenhang zwischen Brute Force Angriffen, den Abwehrmaßnahmen auf ebendiese sowie eventueller Nichtverfügbarkeit des Dienstes. Näheres dazu im Abschnit~\ref{sec:bruteforce}~\nameref{sec:bruteforce}. 

\subsection{Man in the Middle}

Beim Man in the Middle Angriff versucht der Angreifer sich unbemerkt als Mittelsmann in die Kommunikation zwischen Kunde und Webshop einzuklinken. Obwohl der Kunde denkt, direkt mit der Anwendung zu kommunizieren, kann jeglicher Datenfluss vom Angreifer mitgelesen, protokolliert und in schlimsten Fall manipuliert werden. Der Angreifer benötigt daher Zugriff auf die Infrastruktur. Dieser kann im tatsächlichen Zwischenschalten von Abhörhardware\footnote{vgl. \cite{stoll}, S. 28} oder im virtuellen Zwischenschalten auf den unteren Protokollebenen der Netzwerkkommunikation erfolgen.\footnote{vgl. \cite{pritchett}, Kapitel 7.8}

Da der Applikationsentwickler keinen Einfluss auf die Infrastruktur des Internets hat, bleibt hier als Gegenmaßnahme lediglich die Wahl von verschlüsselten Kommunikaitonswegen, z.B. per \ac{HTTPS}.

\subsection{Brute Force}
\label{sec:bruteforce}

Bei einem Brute Force Angriff setzt der Angreifer nicht wie es der Name vermuten lässt auf rohe Gewalt, sondern benötigt Geduld und Ausdauer. Durch Ausprobieren aller möglichen Kombinationen von Benutzername und Passwort versucht der Angreifer Zugriff auf die geschützten Bereiche des Webshops zu erlangen. Es werden oft Listen mit beliebten Passwortern verwendet, um die Effizienz des Angriffs zu optimieren.

Hierbei ist beachtenswert, dass durch die wiederholten Anmeldeversuche der Server überlastet werden kann. Somit kann ein nicht erfolgreicher Brute Force Angriff durchaus zu einem unbeabsichtigten, aber erfolgreichen \ac{DoS}--Angriff entwickeln.\footnote{vgl. \cite{defat}}

Als Gegenmaßnahmen kann die Anzahl der Anmeldeversuche pro Benutzer oder von einer IP--Adresse aus beschränkt werden. Da der Angreifer jedoch meist nicht nur einen, sondern mehrere Nutzernamen durchprobiert, kann hier recht schnell für mehrere Nutzer der Zugriff gesperrt werden. Die regulären Kunden können dann den Webshop nicht erreichen. Die Auswirkungen entsprechen wieder \ac{DoS}.

Der Angreifer kann durch starke Passworte aufgehalten werden. Somit sollte beim Anwendungsentwurf darauf geachtet werden, dass lange Passworte möglich sind, die sowohl Buchstaben und Zahlen als auch Sonderzeichen enthalten können. Je ungewöhnlicher ein Passwort ist, desto wahrscheinlicher ist es, dass es nicht in den Passwortlisten der Angreifer auftaucht, was einen Angriff langsamer, teurer und somit unattraktiver macht. Allerdings sind starke Passworte nicht benutzerfreundlich. Kryptographisch ausreichende Passworte benötigen 20 bis 32 Zeichen und sind somit im Alltag nicht gebräuchlich.\footnote{vgl. \cite{kaufman}, Kapitel 10.4}

Somit ist eine Erweiterung des Passwortbegriffs notwendig, um Brute Force Attacken wirksam bekämpfen zu können. Die Benutzerauthentifikation kann auf zwei oder mehr Stufen erweitert werden. Hierzu wird bei der Benutzeranmeldung nicht nur das abgefragt, was der Nutzer \textit{weiß} (also ein Passwort oder eine Passphrase), sondern auch etwas, was er \textit{besitzt}, was er\textit{ist}\footnote{vgl. \cite{kaufman}, Kapitel 10} oder was er \textit{kann}.

\begin{itemize}
	
\item \textbf{Können:} Eine Abwehrmethode sind \acs{CAPTCHA}s\footnote{\ac{CAPTCHA}}, also für Menschen einfache Aufgaben, die für Computerprogramme sehr schwer zu lösen sind.\footnote{Diese können aus umgangssprachlich Formulierten Rechenaufgaben auf Erstklässlerniveau oder aus dem Erkennen einer Zeichenfolge in einer Grafik mit Hintergrundrauschen bestehen.} \ac{CAPTCHA}s sind schon weit verbreitet, sind aber oft nicht benutzerfreundlich.
\item \textbf{Besitzen:} Schlüssel, Magnetkarten oder Transponder können zur Benutzeranmeldung eingesetzt werden. Desweiteren besteht die Möglichkeit, dass ein Codegenerator\footnote{z.B. Google--Authenticator als Smartphone--App oder spezielle Hardwaregeneratoren} einen Code erzeugt, der nur eine bestimmte, kurze Zeit gültig ist. Somit sind ausgespäte Codes für einen Angreifer wertlos. Nur der Besitzer des Generators kann diese Zeitnah erzeugen und sich damit authentifizieren. 
\item \textbf{Sein:} Über biometrische Merkmale wird die Identität des Benutzers festgestellt.

\end{itemize}

Auch wenn jede dieser Methoden Nachteile hat und für sich alleine genommen nicht zum Aufbau eines Authentifizierungssystems ausreicht, so können diese Methoden zu einer sogenannten Anmeldung in zwei oder mehr Schritten kombiniert werden, und somit eine wirkungsvolle Maßnahme gegen Brute Force Attacken sowie einer Reihe anderer Risiken rund um Passworte darstellen.

\subsection{SQL--Injection}

\subsection{Cross Site Scripting}


\section{Fazit \& Ausblick}

\subsection{Fazit}

Bei Betrachtung der Angriffsmöglichkeiten und der Risiken wird offensichtlich, dass es nur mit beachtlichem Aufwand möglich ist, einen Webshop vollständig vor Angreifern zu schützen. Speziell aufgrund der Tatsache, dass hier nur eine Auswahl der Bedrohungen besprochen wurde.

Andererseits ist auch zu erkennen, dass viele Angriffsszenarien durch die Auswahl entsprechender Komponenten und Best--Practices wirkungsvoll abgewehrt werden können. Wenn Sicherheitsaspekte schon von Beginn an in den Softwareentwurf eingehen, ist die resultierende Applikation zumindest vor den meisten Angriffen geschüzt. Zumindest zufälligen Angreifern, die im Internet nach offensichtlichen Sicherheitslücken suchen, wird somit kein Ziel geboten.

\subsection{Ausblick}

Auf dem Weg zu einer tatsächlichen Implementierung ist es notwendig, weitere Angriffsmethoden zu betrachten, damit das fertige Produkt zumindest gegen die von der OWASP erstellten Top--10 Liste bestehen kann. Die jeweils notwendigen Gegenmaßnahmen sollten als Richtlinien für die Codierung festgelegt und überprüft werden. 

Falls der Internetshop um eine Bezahlfunktion erweitert werden soll, sind noch weitere Angriffsziele zu beachten, da in diesem Fall nicht nur Daten, sondern auch Geld\footnote{Das Geld unserer Kunden} sowie Waren\footnote{Die angebotenen Waren ohne Bezahlung} erbeutet werden können. Dieses Szenario würde aber den Umfang dieser Arbeit überschreiten und müsste somit seperat erörtert werden.

Für weitere Untersuchungen würde sich die Fragestellungen anbieten, in wie weit die am Markt erhältliche Shopsoftware gegen Angriffe geschützt ist bzw. wie groß der Entwicklungsaufwand ist, den am Markt üblichen Grad an Sicherheit bei einer selbst erstellten Anwendung zu erreichen. Hierdurch könnte eine Make-or-Buy Entscheidung unterstützt bzw. begründet werden.
\include{anhang}

\end{spacing}

\clearpage

% Literaturverzeichniss - Ab hier wieder Roemische Seitenzahlen
\pagestyle{plain}
\pagenumbering{roman}
\setcounter{page}{\theromanPagenumber}
\bibliographystyle{apalike}
\bibliography{literatur}
\onehalfspacing
\clearpage

\pagestyle{empty} 
\thispagestyle{empty}

\begin{center}
{\Large Eidesstattliche Erklärung}
\vspace*{4cm}\end{center}
\noindent
Ich versichere, dass ich das beiliegende Assignment selbstständig verfasst, keine anderen als die angegebenen Quellen und Hilfsmittel benutzt sowie alle wörtlich oder sinngemäß übernommenen Stellen in der Arbeit gekennzeichnet habe. 
\vspace{3cm}

\hspace{-0.8cm}
\rule[0.5ex]{6.5cm}{1pt}
\hspace{1.3cm}
\rule[0.5ex]{6.5cm}{1pt}
\\(Datum, Ort)
\hspace{6.3cm}
(Unterschrift)

\clearpage

%Messbox zur Druckkontrolle:
\newcommand{\Messbox}[2]{% Parameters: #1=Breite, #2=Hoehe
\setlength{\unitlength}{1.0mm}%
\begin{picture}(#1,#2)%
\linethickness{0.05mm}%
\put(0,0){\dashbox{0.2}(#1,#2)%
{\parbox{#1mm}{%
\centering\footnotesize 
%{\bf MESSBOX}\\ 
Breite $ = #1 {\rm\ mm}$\\
H\"ohe $ = #2 {\rm\ mm}$
}}}\end{picture}
}

\begin{center}
\textbf{--- Druckgröße kontrollieren! ---}
\\
\Messbox{100}{50} % Angabe der Breite/Hoehe in mm
\\
\textbf{--- Diese Seite nach dem Druck entfernen! ---}
\end{center}


\end{document}

